\subsection{Pembagian Fungsi}
\begin{definition}[Pembagian Fungsi]
  \[ \left(\frac{f}{g}\right)(x) = \frac{f(x)}{g(x)} \]
  Dengan syarat $g(x) \ne 0$.
\end{definition}

Ada 3 langkah mengerjakan pembagian fungsi aljabar yaitu:
\begin{enumerate}
\item Menuliskan nilai $\frac{f(x)}{g(x)}$.
\item Menyederhanakan pecahan (dengan memfaktorkan).
\item Membagi dan menuliskan hasil akhir.
\end{enumerate}

\subsubsection{Cara Memfaktorkan}
\begin{enumerate}
\item \textbf{Ambil angka/variabel yang sama dari semua suku.}
  \begin{example}[Variabel yang Sama]
    \hfill 
    \begin{enumerate}
    \item $6x + 9 = 3(2x + 3)$
    \item $4x^2 - 8x = 4x(x - 2)$
    \item $x^3 + 2x^2 = x^2(x + 2)$
    \end{enumerate}
  \end{example}

\item \textbf{Selisih 2 kuadrat.}
  \[ a^2 - b^2 = (a + b)(a - b) \]
  \begin{example}[Selisih 2 Kuadrat]
    \hfill
    \begin{enumerate}
    \item $x^2 - 9 = (x + 3)(x - 3)$
    \item $x^2 - 16 = (x + 4)(x - 4)$
    \item $9x^2 - 25 = (3x + 5)(3x - 5)$
    \end{enumerate}
  \end{example}

\item \textbf{Trinomial $ax^2 + bx + c = 0 \quad (a = 1)$.} \\
  Cari 2 bilangan yang jika dikalikan hasilnya $c$ dan jika dijumlahkan hasilnya $b$.
  \begin{example}[Trinomial]
    \hfill \\
    Faktorkan $x^2 + 5x + 6 = 0$:
    \[
    \begin{aligned}
      6 \cdot 1 = 6 \quad&\implies\quad 6 + 1 = 7 \quad \text{(Salah)} \\
      3 \cdot 2 = 6 \quad&\implies\quad 3 + 2 = 5 \quad \text{(Benar)}
    \end{aligned}
    \]
    Sehingga bentuk faktornya adalah: $(x + 3)(x + 2)$.
  \end{example}

\item \textbf{Trinomial $a \ne 1$.}
  \begin{enumerate}
  \item Pertama hitung $a \cdot c$.
  \item Cari 2 bilangan yang dikalikan hasilnya $ac$ dan dijumlahkan hasilnya $b$.
  \item Pecahkan suku tengah menjadi 2 suku, lalu kelompokkan.
  \end{enumerate}
  \begin{example}[Trinomial $a \ne 1$]
    \hfill \\
    Sederhanakan $\left(\frac{f}{g}\right)(x) = \frac{6x^2 + x - 2}{2x - 1}$:
    \[
    \begin{aligned}
      \frac{6x^2 + x - 2}{2x - 1} &= \frac{6x^2 - 3x + 4x - 2}{2x - 1} \\
      &= \frac{3x(2x - 1) + 2(2x - 1)}{2x - 1} \\
      &= \frac{(3x + 2)(2x - 1)}{2x - 1} \\
      &= 3x + 2
    \end{aligned}
    \]
  \end{example}
\end{enumerate}
